\section{绪论 20250922}

\subsection{建立与发展}

\begin{enumerate}
    \item 从宏观到微观；
    \item 从体相到表相；
    \item 从静态到动态；
    \item 从定性到定量；
    \item 从单一学科到边缘交叉学科；
    \item 从平衡态的研究到非平衡态的研究。
\end{enumerate}

\subsection{目的与内容}

\begin{enumerate}
    \item 用物理的理论和实验方法研究化学变化的本质和规律；
    \item 化学变化的方向和限度问题；
    \item 化学反应的速率和机理问题；
    \item 物质结构和性能之间的关系。
\end{enumerate}

\subsection{研究方法}

\begin{enumerate}
    \item 归纳法、演绎法、$\cdots$
\end{enumerate}

\subsection{学习方法}

\begin{enumerate}
    \item 扩大知识面，打好专业基础；
    \item 提高自学能力，培养独立思考、观察、联想能力；
    \item 抓住每章重点；
    \item 掌握主要公式的物理意义和使用条件；
    \item 课前自学，做笔记，及时复习；
    \item 注意章节间联系，融会贯通；
    \item 重视习题，培养独立思考能力。
\end{enumerate}
%----------%